\section{重み三への特殊化と幾何的辞書}
\label{sec:poisson-coordinates-v32}
以下の成分分類では $m=3$ とし、$X=X_3$、$\mathscr P=\mathscr P_3$ と書く。
命題~\ref{prop:general-equations-v33}は、五十三変数の標準形
\[
 c\in S_3,\quad D\in\mathfrak{sl}_3,\quad [F]\in S_4^{\oplus3}/S_3\mathbf x
\]
と三十八個の独立な二次式を与える。
\begin{proposition}[重み三の方程式]
\label{prop:weighted-poisson-equations-v31}
全Jacobi式は
\begin{align}
 D(c)&=0,\label{eq:weighted-base-jacobi-v31}\\
 2F(c)-(\mathbf x\times D)\cdot F-c\operatorname{div}F&=0.
 \label{eq:weighted-mixed-jacobi-v31}
\end{align}
係数ベクトルは $p=(\nabla c-\mathbf x\times D)/3$ である。
\end{proposition}
証明は一般重みの命題で与えた。
$S_2^{\oplus3}\simeq S_3\oplus\mathfrak{sl}_3$ という三変数二次Poisson構造の分解は
\cite{LiuXu1992,DoninMakarLimanov1998}と関係する。
